\documentclass[11pt,oneside]{article}

\usepackage[a4paper,margin=27mm,headheight=15pt]{geometry}
\usepackage[T1]{fontenc}
\usepackage[utf8]{inputenc}
\IfFileExists{libertinus.sty}{\usepackage{libertinus}}{\usepackage{lmodern}}
\usepackage{microtype}
\usepackage{amsmath,amssymb,amsthm,mathtools,mathrsfs}
\usepackage{bm}
\usepackage{booktabs,longtable,array,tabularx}
\usepackage{xcolor}
\usepackage{enumitem}
\usepackage{fancyhdr}
\usepackage{titlesec}
\usepackage{listings}
\usepackage{xurl}
\usepackage{hyperref}
\usepackage[nameinlink,noabbrev]{cleveref}

\definecolor{AtlasBlue}{RGB}{31,72,112}
\definecolor{AtlasGray}{RGB}{90,96,104}
\definecolor{AtlasLight}{RGB}{241,245,248}
\hypersetup{
  colorlinks=true,
  linkcolor=AtlasBlue,
  citecolor=AtlasBlue,
  urlcolor=AtlasBlue,
  pdftitle={A Formula Atlas for the Thue--Morse Sequence},
  pdfauthor={Vladimir Reshetnikov / consolidated research draft}
}
\setlist{nosep,leftmargin=2em}
\setlength{\parindent}{0pt}
\setlength{\parskip}{0.55em}
\allowdisplaybreaks
\numberwithin{equation}{section}

\pagestyle{fancy}
\fancyhf{}
\fancyhead[L]{\small A Formula Atlas for the Thue--Morse Sequence}
\fancyhead[R]{\small\thepage}
\renewcommand{\headrulewidth}{0.35pt}
\titleformat{\section}{\Large\bfseries\color{AtlasBlue}}{\thesection}{0.7em}{}
\titleformat{\subsection}{\large\bfseries\color{AtlasBlue}}{\thesubsection}{0.7em}{}

\newtheorem{theorem}{Theorem}[section]
\newtheorem{proposition}[theorem]{Proposition}
\newtheorem{lemma}[theorem]{Lemma}
\newtheorem{corollary}[theorem]{Corollary}
\newtheorem{identity}[theorem]{Identity}
\theoremstyle{definition}
\newtheorem{definition}[theorem]{Definition}
\newtheorem{example}[theorem]{Example}
\newtheorem{algorithm}[theorem]{Algorithm}
\theoremstyle{remark}
\newtheorem{remark}[theorem]{Remark}
\newtheorem{literature}[theorem]{Literature theorem}

\newcommand{\N}{\mathbb N}
\newcommand{\Z}{\mathbb Z}
\newcommand{\Q}{\mathbb Q}
\newcommand{\R}{\mathbb R}
\newcommand{\C}{\mathbb C}
\newcommand{\F}{\mathbb F}
\newcommand{\wt}{\operatorname{wt}}
\newcommand{\vtwo}{\nu_2}
\newcommand{\eps}{\varepsilon}
\newcommand{\TM}{\tau}
\newcommand{\ii}{\mathrm i}
\newcommand{\e}{\mathrm e}
\newcommand{\dd}{\,\mathrm d}
\newcommand{\ind}{\mathbf 1}
\newcommand{\xor}{\mathbin{\mathsf{xor}}}
\newcommand{\bitand}{\mathbin{\mathsf{and}}}
\newcommand{\submask}{\preceq_2}
\newcommand{\Bell}{\mathcal B}
\newcommand{\fall}[2]{#1^{\underline{#2}}}
\newcommand{\bigO}{\mathcal O}
\newcommand{\smallo}{o}
\newcommand{\resmod}[2]{[\,#1\,]_{#2}}
\newcommand{\sinc}{\operatorname{sinc}}
\newcommand{\repo}{\textsc{ProveIt}}
\newcommand{\lean}[1]{\texttt{#1}}
\newcommand{\newmaterial}{\textcolor{AtlasBlue}{\textbf{Added in this synthesis.}}\ }

\lstset{
  basicstyle=\ttfamily\small,
  columns=fullflexible,
  frame=single,
  backgroundcolor=\color{AtlasLight},
  breaklines=true,
  showstringspaces=false
}

\begin{document}
\pagenumbering{Roman}

\begin{titlepage}
\centering
\vspace*{2.2cm}
{\Huge\bfseries\color{AtlasBlue} A Formula Atlas for the\\[0.25em] Thue--Morse Sequence\par}
\vspace{1cm}
{\Large Exact arithmetic, Pascal-triangle, product, Prouhet,\\
Fourier, spectral, infinite-product, and Fabius representations\par}
\vspace{1.5cm}
{\large A consolidated and expanded version of four research drafts\par}
\vspace{0.4cm}
{\large in Vladimir Reshetnikov's \repo\ project\par}
\vfill
\begin{minipage}{0.82\textwidth}
\small
This manuscript merges the four files \emph{Further Formulae for the
Thue--Morse Sequence--1} through \emph{--4}.  Repeated statements have been
collapsed, notation has been standardized, proofs have been cross-checked,
and the material has been arranged as a dependency-respecting atlas rather
than as four parallel reports.  New material includes a minimal-kernel
viewpoint, a Catalan-valuation formula, a direct natural-boundary proof,
entire shifted Dirichlet series, additional Boolean-lattice inversions,
word-complexity consequences, a base-
\(B\) root-of-unity Prouhet theorem, and a modern distribution survey.
No claim of historical priority is intended.
\end{minipage}
\vfill
{\large 26 August 2026\par}
\end{titlepage}
\pagenumbering{roman}

\begin{abstract}
Let \(\TM(n)=\wt(n)\bmod 2\) be the zero--one Thue--Morse bit and
\(\eps_n=(-1)^{\wt(n)}=1-2\TM(n)\) its signed form.  Almost every useful
formula becomes transparent after moving among five equivalent objects:
binary digit parity, a character of the Boolean cube, the finite product
\(\prod_{j<m}(1-z^{2^j})\), the substitution \(0\mapsto01,1\mapsto10\),
and a mixed finite-difference operator.  The purpose of this document is to
make those translations explicit and reusable.  It develops arithmetic and
Pascal identities, exact prefix discrepancy, valuation recurrences, partition
and determinant formulas, Prouhet cancellation and all-order moments,
ordinary and Walsh Fourier transforms, Riesz products and autocorrelation,
Mahler functions, Dirichlet and product transforms, finite-state and word
combinatorics, base-\(B\) extensions, and the discrete bridge to the Fabius
function.  Results imported from the literature are labelled as such; all
other displayed identities are proved or reduced to a previously proved
identity in the text.
\end{abstract}

\tableofcontents
\clearpage
\pagenumbering{arabic}

\section{Conventions, provenance, and a one-page map}
\label{sec:conventions}

Throughout \(\N=\{0,1,2,\ldots\}\).  If
\[
 n=\sum_{j\ge0}b_j(n)2^j,\qquad b_j(n)\in\{0,1\},
\]
then
\begin{equation}
 \wt(n)=\sum_{j\ge0}b_j(n),\qquad
 \TM(n)=\wt(n)\bmod2,\qquad
 \eps_n=(-1)^{\wt(n)}=1-2\TM(n).
 \label{eq:definitions}
\end{equation}
The terminology \emph{evil} and \emph{odious} refers respectively to
\(\TM(n)=0\) and \(\TM(n)=1\).

The four source drafts were independently organized and consequently overlap
in their basic recurrences, Lucas/Kummer formulas, finite products, Prouhet
identities, Fourier products, prefix sums, and Woods--Robbins evaluation.  The
present text retains one proof of each central theorem and places specialized
consequences immediately after the theorem that generates them.  The source
files are available at
\begin{center}
\url{https://github.com/VladimirReshetnikov/ProveIt/tree/main/Analysis/FabiusFunction/docs/non-formalized-research-frontiers/drafts}.
\end{center}

\subsection{Formula map}

\begin{longtable}{@{}p{0.20\textwidth}p{0.33\textwidth}p{0.38\textwidth}@{}}
\toprule
Viewpoint & Exact representation & Main use \\
\midrule
\endhead
Binary recursion & \(\eps_{2n}=\eps_n,\ \eps_{2n+1}=-\eps_n\) & Fast computation, blocks, automata \\
Boolean character & \(\eps_n=(-1)^{\bm1\cdot\bm b(n)}\) & XOR, Walsh transform, submasks \\
Valuation & \(\wt(n)=\vtwo\binom{2n}{n}\) & Binomial/Catalan formulas \\
Pascal row & \(\#\{k:\binom nk\text{ odd}\}=2^{\wt(n)}\) & Binomial-log and log-free formulas \\
Finite product & \(P_m(z)=\prod_{j<m}(1-z^{2^j})\) & Moments, DFT, Riesz products \\
Infinite product & \(E(z)=\prod_{j\ge0}(1-z^{2^j})\) & Mahler equation, constants, transforms \\
Finite difference & \(\sum_{n<2^m}\eps_nf(n)=(-1)^m\Delta_1\cdots\Delta_{2^{m-1}}f(0)\) & Prouhet, remainders, reciprocal sums \\
Prefix transform & \(\sum_nS_n^{(k)}z^n=E(z)/(1-z)^k\) & Discrete antiderivatives, Fabius bridge \\
Spectral measure & \(2^{-m}|P_m(e^{2\pi ix})|^2\dd x\) & Autocorrelation, singular spectrum \\
\bottomrule
\end{longtable}

\begin{remark}[Proof status]
A label ``Literature theorem'' means that only context or a proof sketch is
given.  ``Added in this synthesis'' means absent from the four source drafts
as a named result; it does not assert novelty in the historical literature.
The Lean file names mentioned near the end are a formalization roadmap, not a
claim that every identity in this atlas is already formalized.
\end{remark}

\section{Binary recursion, blocks, automata, and uniqueness}
\label{sec:binary}

\subsection{The two-scale recurrence and substitution}

\begin{theorem}[Fundamental recursion]
For every \(n\ge0\),
\begin{equation}
 \boxed{\eps_{2n}=\eps_n,\qquad \eps_{2n+1}=-\eps_n,}
 \label{eq:eps-recursion}
\end{equation}
and equivalently
\begin{equation}
 \TM(2n)=\TM(n),\qquad \TM(2n+1)=1-\TM(n).
 \label{eq:tau-recursion}
\end{equation}
\end{theorem}

\begin{proof}
Appending a binary zero does not change the digit sum, while appending a
binary one increases it by one.  Taking parity gives
\eqref{eq:eps-recursion} and \eqref{eq:tau-recursion}.
\end{proof}

Starting from \(0\), the morphism
\begin{equation}
 \mu(0)=01,\qquad \mu(1)=10
 \label{eq:morphism}
\end{equation}
produces the zero--one word.  In signed notation the corresponding operation
replaces a block \(W\) by \(W,-W\).  Thus the block of length \(2^{m+1}\) is
the block of length \(2^m\) followed by its complement.

\subsection{Binary concatenation}

\begin{theorem}[Block multiplicativity]
If \(m,q\ge0\) and \(0\le r<2^m\), then
\begin{equation}
 \boxed{\eps_{2^mq+r}=\eps_q\eps_r,}
 \qquad
 \TM(2^mq+r)\equiv \TM(q)+\TM(r)\pmod2.
 \label{eq:block-multiplicativity}
\end{equation}
\end{theorem}

\begin{proof}
The binary digits of \(2^mq+r\) are the digits of \(q\), shifted by \(m\)
places, followed by the \(m\)-digit expansion of \(r\); the supports are
disjoint, so their weights add.
\end{proof}

In particular,
\begin{equation}
 (\eps_{2^m},\ldots,\eps_{2^{m+1}-1})
 =-(\eps_0,\ldots,\eps_{2^m-1}).
 \label{eq:complement-block}
\end{equation}

\subsection{The exact two-kernel and a minimal automaton}

\begin{theorem}[Exact two-kernel]
For \(e\ge0\) and \(0\le r<2^e\),
\begin{equation}
 \eps_{2^en+r}=\eps_r\eps_n.
 \label{eq:kernel-formula}
\end{equation}
Consequently the \(2\)-kernel of the signed sequence consists of exactly the
two sequences \((\eps_n)_{n\ge0}\) and \((-\eps_n)_{n\ge0}\).
\end{theorem}

\begin{proof}
Apply \cref{eq:block-multiplicativity}.  The sign is \(\eps_r\in\{\pm1\}\).
Both possibilities occur, for example at \(r=0\) and \(r=1\).
\end{proof}

\newmaterial The preceding theorem makes the minimal automaton completely
explicit.  Let
\[
 M_0=\begin{pmatrix}1&0\\0&1\end{pmatrix},\qquad
 M_1=\begin{pmatrix}0&1\\1&0\end{pmatrix},\qquad e_0=\binom10.
\]
If \(n=\sum_{j=0}^{m-1}b_j2^j\), then
\begin{equation}
 v(n)=M_{b_0}\cdots M_{b_{m-1}}e_0
 =\binom{1-\TM(n)}{\TM(n)},
 \label{eq:automaton-vector}
\end{equation}
so
\begin{equation}
 \TM(n)=\begin{pmatrix}0&1\end{pmatrix}v(n),\qquad
 \eps_n=\begin{pmatrix}1&-1\end{pmatrix}v(n).
 \label{eq:automaton-output}
\end{equation}
The two states record only the parity already read; input \(1\) toggles the
state and input \(0\) fixes it.  Two states are necessary because the outputs
at \(0\) and \(1\) differ.

\subsection{Nonperiodicity}

\begin{theorem}[No eventual period]
The Thue--Morse sequence is not eventually periodic.
\end{theorem}

\begin{proof}
Suppose that \(p>0\) were an eventual period.  Fix \(q\ge0\) and choose
\(k\) so large that \(2^k+q\) and \(2^k+q+p\) lie in the periodic tail and
\(q+p<2^k\).  Block multiplicativity gives
\[
 -\eps_q=\eps_{2^k+q}=\eps_{2^k+q+p}=-\eps_{q+p},
\]
so \(p\) would in fact be a period from the beginning.  Choose \(k\) with
\(2^k>p\).  Periodicity would imply
\(\eps_{2^k-p}=\eps_{2^k}=-1\).  Yet
\(2^k-p=(2^k-1)-(p-1)\), so its \(k\) digits are the complements of the
\(k\) digits of \(p-1\), and
\[
 \eps_{2^k-p}=(-1)^{k-\wt(p-1)}.
\]
This alternates when \(k\) is increased by one, a contradiction.
\end{proof}

\section{Floor, fractional-part, trigonometric, and bitwise formulas}
\label{sec:digits}

\subsection{Digit extraction}

For every \(j\ge0\),
\begin{equation}
 b_j(n)=\left\lfloor\frac{n}{2^j}\right\rfloor
       -2\left\lfloor\frac{n}{2^{j+1}}\right\rfloor.
 \label{eq:digit-floor}
\end{equation}
Therefore, with all but finitely many factors equal to \(1\),
\begin{equation}
 \boxed{\eps_n=
 \prod_{j\ge0}\left(1-2\left\lfloor\frac{n}{2^j}\right\rfloor
 +4\left\lfloor\frac{n}{2^{j+1}}\right\rfloor\right).}
 \label{eq:pure-floor-product}
\end{equation}
Also,
\begin{equation}
 \wt(n)=n-\sum_{j\ge1}\left\lfloor\frac n{2^j}\right\rfloor
       =\sum_{j\ge1}\left\{\frac n{2^j}\right\},
 \label{eq:weight-floor-frac}
\end{equation}
where \(\{x\}=x-\lfloor x\rfloor\).  The second equality follows because
\(\sum_{j\ge1}n/2^j=n\).

Consequently,
\begin{equation}
 \eps_n=(-1)^{\,n+\sum_{j\ge1}\lfloor n/2^j\rfloor}
 =\cos\!\left(\pi\sum_{j\ge1}\left\{\frac n{2^j}\right\}\right),
 \qquad
 \TM(n)=\frac{1-\eps_n}{2}.
 \label{eq:floor-trig}
\end{equation}

\subsection{A sign product and an entire sinc product}

\begin{proposition}[Rademacher--sine sign product]
For every \(n\ge0\),
\begin{equation}
 \boxed{\eps_n=
  \prod_{j\ge0}\operatorname{sgn}
  \sin\!\left(\frac{\pi(n+\tfrac12)}{2^j}\right).}
 \label{eq:sine-sign-product}
\end{equation}
Only finitely many factors are negative.
\end{proposition}

\begin{proof}
The sign of the \(j\)-th sine is
\((-1)^{\lfloor(n+1/2)/2^j\rfloor}=(-1)^{\lfloor n/2^j\rfloor}\).
Multiplying and using \eqref{eq:weight-floor-frac} modulo two gives
\((-1)^{\wt(n)}\).
\end{proof}

\newmaterial Define \(\sinc y=\sin y/y\), with \(\sinc0=1\), and
\begin{equation}
 \Phi(x)=\prod_{j\ge0}\sinc\!\left(\frac{\pi x}{2^j}\right).
 \label{eq:Phi-sinc}
\end{equation}
The product converges locally uniformly and hence defines an entire function.
For every integer \(N\ge0\) and \(N<x<N+1\),
\begin{equation}
 \boxed{\operatorname{sgn}\Phi(x)=\eps_N.}
 \label{eq:sinc-sign}
\end{equation}
Indeed, the denominator of each factor is positive and the signs of the
numerators multiply exactly as in \eqref{eq:sine-sign-product}, with \(x\)
in place of \(N+1/2\).  Thus the Thue--Morse word is the sign itinerary of a
single canonical entire product between its successive integer zeros.

\subsection{XOR and carry-free multiplication}

For nonnegative integers \(a,b\),
\begin{equation}
 \wt(a\xor b)=\wt(a)+\wt(b)-2\wt(a\bitand b).
 \label{eq:xor-weight}
\end{equation}
Hence
\begin{equation}
 \boxed{\eps_{a\xor b}=\eps_a\eps_b,\qquad
 \TM(a\xor b)\equiv\TM(a)+\TM(b)\pmod2.}
 \label{eq:xor-character}
\end{equation}
This says that \(n\mapsto\eps_n\), restricted to \(m\)-bit integers, is the
Walsh character indexed by the all-ones vector.

If \(a\bitand b=0\), then \(a+b=a\xor b\), and therefore
\begin{equation}
 \eps_{a+b}=\eps_a\eps_b.
 \label{eq:carry-free-addition}
\end{equation}
The correction in the presence of carries will appear in
\cref{sec:valuations}.

\section{Two-adic valuations, carries, and Catalan numbers}
\label{sec:valuations}

\subsection{Legendre and the central binomial coefficient}

Legendre's formula in base two is
\begin{equation}
 \vtwo(n!)=\sum_{j\ge1}\left\lfloor\frac n{2^j}\right\rfloor
          =n-\wt(n).
 \label{eq:legendre}
\end{equation}
It immediately yields one of the cleanest arithmetic encodings.

\begin{theorem}[Central-binomial valuation]
For every \(n\ge0\),
\begin{equation}
 \boxed{\vtwo\binom{2n}{n}=\wt(n),\qquad
 \eps_n=(-1)^{\vtwo\binom{2n}{n}}.}
 \label{eq:central-binomial}
\end{equation}
\end{theorem}

\begin{proof}
By \eqref{eq:legendre},
\[
 \vtwo\binom{2n}{n}
 =(2n-\wt(2n))-2(n-\wt(n))=\wt(n),
\]
because \(\wt(2n)=\wt(n)\).
\end{proof}

\subsection{Catalan valuation}

\begin{theorem}[Catalan-number formula]
\newmaterial Let \(C_n=\frac1{n+1}\binom{2n}{n}\).  Then
\begin{equation}
 \boxed{\vtwo(C_n)=\wt(n+1)-1
 =\wt(n)-\vtwo(n+1).}
 \label{eq:catalan-valuation}
\end{equation}
Consequently
\begin{equation}
 \eps_{n+1}=-(-1)^{\vtwo(C_n)},
 \qquad
 \eps_n=(-1)^{\vtwo(C_n)+\vtwo(n+1)}.
 \label{eq:catalan-thue-morse}
\end{equation}
\end{theorem}

\begin{proof}
Use \eqref{eq:central-binomial} and subtract \(\vtwo(n+1)\).  When one is
added to \(n\), its trailing block of \(\vtwo(n+1)\) one-bits becomes zero
and the next zero becomes one, so
\(\wt(n+1)=\wt(n)+1-\vtwo(n+1)\).
\end{proof}

\subsection{Kummer carries and multinomials}

Kummer's theorem gives
\begin{equation}
 \vtwo\binom{a+b}{a}=\wt(a)+\wt(b)-\wt(a+b),
 \label{eq:kummer-two}
\end{equation}
which is the number of binary carries, counted with multiplicity through
propagation.  Taking parity yields the exact carry correction
\begin{equation}
 \boxed{\eps_{a+b}=\eps_a\eps_b
 (-1)^{\vtwo\binom{a+b}{a}}.}
 \label{eq:carry-correction}
\end{equation}
More generally, if \(k_1+\cdots+k_r=n\), then
\begin{equation}
 \vtwo\binom{n}{k_1,\ldots,k_r}
 =\sum_{i=1}^r\wt(k_i)-\wt(n),
 \label{eq:multinomial-valuation}
\end{equation}
and therefore
\begin{equation}
 \boxed{\eps_n=
  \left(\prod_{i=1}^r\eps_{k_i}\right)
  (-1)^{\vtwo\binom{n}{k_1,\ldots,k_r}}.}
 \label{eq:multinomial-sign}
\end{equation}

\section{Pascal-triangle formulas and Boolean-lattice inversion}
\label{sec:pascal}

\subsection{Odd entries and the binomial--logarithm formula}

Lucas's theorem modulo two says
\begin{equation}
 \binom nk\equiv1\pmod2
 \quad\Longleftrightarrow\quad k\submask n,
 \label{eq:lucas-submask}
\end{equation}
meaning that each one-bit of \(k\) is also a one-bit of \(n\).  Hence the
number of odd entries in row \(n\) of Pascal's triangle is
\begin{equation}
 \boxed{\#\left\{0\le k\le n:\binom nk\text{ is odd}\right\}
 =2^{\wt(n)}.}
 \label{eq:odd-row-count}
\end{equation}

Set
\begin{equation}
 X(n)=\sum_{k=0}^n(-1)^{\binom nk}.
 \label{eq:X-def}
\end{equation}
Every even entry contributes \(+1\), every odd entry \(-1\), so
\begin{equation}
 \boxed{X(n)=n+1-2^{\wt(n)+1}.}
 \label{eq:X-exact}
\end{equation}
This recovers the source document's binomial--logarithm formula:
\begin{equation}
 \TM(n)=\frac{1+(-1)^{\log_2(n+1-X(n))}}2,
 \label{eq:binomial-log}
\end{equation}
because \(n+1-X(n)=2^{\wt(n)+1}\).  The logarithm is exact, not numerical.

\subsection{A log-free residue formula}

Let
\begin{equation}
 S(n)=1+X(n)=n+2-2^{\wt(n)+1}.
 \label{eq:S-row}
\end{equation}
Since powers of two alternate between \(1\) and \(2\) modulo three,
\begin{equation}
 2n+S(n)\equiv
 \begin{cases}0\pmod3,&\wt(n)\text{ even},\\
               1\pmod3,&\wt(n)\text{ odd}.
 \end{cases}
\end{equation}
Thus the least nonnegative residue already is a bit:
\begin{equation}
 \boxed{\TM(n)=\resmod{2n+1+\sum_{k=0}^n(-1)^{\binom nk}}{3}.}
 \label{eq:log-free-mod3}
\end{equation}
Equivalently, because the residue can only be \(0\) or \(1\),
\begin{equation}
 \TM(n)=\frac2{\sqrt3}
 \sin\!\left(\frac{2\pi}{3}\left(2n+1+
 \sum_{k=0}^n(-1)^{\binom nk}\right)\right).
 \label{eq:log-free-sine}
\end{equation}

\subsection{Row enumerators, powers-of-two columns, and submasks}

For an indeterminate \(u\), define the parity enumerator
\[
 R_n(u)=\sum_{k=0}^n u^{\binom nk\bmod2}.
\]
Then
\begin{equation}
 \boxed{R_n(u)=n+1+(u-1)2^{\wt(n)}.}
 \label{eq:row-enumerator}
\end{equation}
Specializing \(u=-1\) gives \eqref{eq:X-exact}.

Lucas's theorem at the power-of-two columns reads
\begin{equation}
 \binom{n}{2^j}\bmod2=b_j(n).
 \label{eq:power-two-column}
\end{equation}
Therefore
\begin{equation}
 \boxed{\eps_n=\prod_{j\ge0}(-1)^{\binom{n}{2^j}}.}
 \label{eq:power-two-column-product}
\end{equation}

The complete submask enumerator is
\begin{equation}
 \boxed{\sum_{k\submask n}z^{\wt(k)}=(1+z)^{\wt(n)}.}
 \label{eq:submask-weight-enumerator}
\end{equation}
At \(z=-2\),
\begin{equation}
 \eps_n=\sum_{k\submask n}(-2)^{\wt(k)}.
 \label{eq:submask-minus-two}
\end{equation}
Combining this with \eqref{eq:central-binomial} and the indicator for oddness
produces a formula involving only binomial coefficients and two-adic
valuations:
\begin{equation}
 \boxed{\eps_n=\sum_{k=0}^n
 \frac{1-(-1)^{\binom nk}}2
 (-2)^{\vtwo\binom{2k}{k}}.}
 \label{eq:pure-binomial-formula}
\end{equation}

\subsection{Pascal parity as a zeta matrix}

Let
\begin{equation}
 P_{n,k}=\ind_{k\submask n}=\binom nk\bmod2.
 \label{eq:pascal-zeta}
\end{equation}
On every finite initial Boolean lattice, \(P\) is the zeta matrix.  Its inverse
is the Möbius matrix
\begin{equation}
 (P^{-1})_{n,k}=P_{n,k}(-1)^{\wt(n)-\wt(k)}.
 \label{eq:pascal-inverse}
\end{equation}
The first column is precisely
\begin{equation}
 (P^{-1})_{n,0}=(-1)^{\wt(n)}=\eps_n.
 \label{eq:pascal-inverse-first-column}
\end{equation}
This places the signed Thue--Morse sequence inside the elementary incidence
algebra of the Boolean lattice.

\subsection{Odd multinomial support}

For fixed \(r\ge2\), an \(r\)-nomial coefficient
\(\binom{n}{k_1,\ldots,k_r}\) is odd exactly when each one-bit of \(n\) is
assigned to exactly one of the \(r\) parts.  Hence
\begin{equation}
 \boxed{\#\left\{(k_1,\ldots,k_r):\sum k_i=n,
 \binom{n}{k_1,\ldots,k_r}\text{ odd}\right\}=r^{\wt(n)}.}
 \label{eq:odd-multinomial-count}
\end{equation}
Since the total number of weak compositions is \(\binom{n+r-1}{r-1}\),
\begin{equation}
 \sum_{k_1+\cdots+k_r=n}
 (-1)^{\binom{n}{k_1,\ldots,k_r}}
 =\binom{n+r-1}{r-1}-2r^{\wt(n)}.
 \label{eq:multinomial-signed-sum}
\end{equation}
For any odd \(r\equiv-1\pmod4\), the parity of \(\wt(n)\) can again be
recovered from a small residue of the right-hand side; the case \(r=3\) is the
multinomial analogue of \eqref{eq:log-free-mod3}.

\section{Exact balance, interval discrepancy, and dyadic blocks}
\label{sec:balance}

Define the half-open signed prefix discrepancy
\begin{equation}
 \mathcal D(N)=\sum_{0\le n<N}\eps_n.
 \label{eq:prefix-D}
\end{equation}

\begin{theorem}[Exact prefix discrepancy]
For every \(q\ge0\),
\begin{equation}
 \boxed{\mathcal D(2q)=0,\qquad
 \mathcal D(2q+1)=\eps_q.}
 \label{eq:prefix-exact}
\end{equation}
In particular \(|\mathcal D(N)|\le1\).
\end{theorem}

\begin{proof}
The first \(2q\) terms cancel in pairs because
\(\eps_{2j}+\eps_{2j+1}=0\).  The next term is
\(\eps_{2q}=\eps_q\).
\end{proof}

Let \(A(N)=\sum_{n<N}\TM(n)\) count odious indices in the prefix.  Since
\(\TM=(1-\eps)/2\),
\begin{equation}
 \boxed{A(2q)=q,\qquad A(2q+1)=q+\TM(q).}
 \label{eq:prefix-one-count}
\end{equation}
Thus every even prefix is exactly balanced.

For \(0\le a\le b\),
\begin{equation}
 \sum_{n=a}^{b-1}\eps_n=\mathcal D(b)-\mathcal D(a),
 \qquad
 \left|\sum_{n=a}^{b-1}\eps_n\right|\le2.
 \label{eq:interval-discrepancy}
\end{equation}
Consequently the number of ones in any interval differs from half the interval
length by at most one.  The signed bound two is sharp, for example on
\([1,3)\).

Block multiplicativity strengthens this to an exact local formula.  If
\(0\le r\le2^m\), then
\begin{equation}
 \boxed{\sum_{j=0}^{r-1}\eps_{2^mq+j}=\eps_q\mathcal D(r).}
 \label{eq:partial-block}
\end{equation}
Every aligned block of length \(2^m\), \(m\ge1\), has zero signed sum, and
every partial block has discrepancy zero or one up to a global sign.

The ordinary generating series of the discrepancy is especially sparse:
\begin{equation}
 \sum_{N\ge0}\mathcal D(N)z^N
 =\frac{z}{1-z}E(z)=zE(z^2)
 =z\prod_{j\ge1}(1-z^{2^j}),
 \qquad |z|<1,
 \label{eq:prefix-generating-D}
\end{equation}
where \(E\) is introduced in \cref{sec:products}.  Coefficient comparison
recovers \eqref{eq:prefix-exact}.

\section{Finite and infinite Euler products}
\label{sec:products}

\subsection{The finite mother product}

For \(m\ge0\), define
\begin{equation}
 P_m(z)=\sum_{n=0}^{2^m-1}\eps_nz^n.
 \label{eq:Pm-def}
\end{equation}

\begin{theorem}[Finite Euler product]
\begin{equation}
 \boxed{P_m(z)=\prod_{j=0}^{m-1}(1-z^{2^j}).}
 \label{eq:finite-product}
\end{equation}
\end{theorem}

\begin{proof}
When the product is expanded, choosing \(-z^{2^j}\) from the factors indexed
by a subset \(J\) produces the unique exponent
\(n=\sum_{j\in J}2^j\) with coefficient \((-1)^{|J|}=\eps_n\).
\end{proof}

The recursion
\begin{equation}
 P_{m+1}(z)=(1-z^{2^m})P_m(z)
 \label{eq:Pm-recursion}
\end{equation}
is the generating-polynomial form of block complementation.  If
\(d=2^m-1\), then
\begin{equation}
 z^dP_m(z^{-1})=(-1)^mP_m(z),
 \qquad
 \eps_{d-n}=(-1)^m\eps_n.
 \label{eq:Pm-reciprocity}
\end{equation}
The cyclotomic factorization is
\begin{equation}
 P_m(z)=(1-z)^m\prod_{r=1}^{m-1}\Phi_{2^r}(z)^{m-r},
 \label{eq:cyclotomic-Pm}
\end{equation}
where \(\Phi_d\) denotes the \(d\)-th cyclotomic polynomial.

\subsection{The infinite product and Mahler equation}

For \(|z|<1\), coefficient stabilization gives
\begin{equation}
 \boxed{E(z)=\sum_{n\ge0}\eps_nz^n
 =\prod_{j\ge0}(1-z^{2^j}).}
 \label{eq:E-product}
\end{equation}
It satisfies the first-order Mahler equation
\begin{equation}
 \boxed{E(z)=(1-z)E(z^2),\qquad E(0)=1.}
 \label{eq:E-mahler}
\end{equation}
The zero--one generating series is
\begin{equation}
 T(z)=\sum_{n\ge0}\TM(n)z^n
 =\frac12\left(\frac1{1-z}-E(z)\right).
 \label{eq:T-from-E}
\end{equation}

\subsection{A direct natural-boundary theorem}

\begin{theorem}[Natural boundary]
\newmaterial Every point of \(|z|=1\) is a singular point of analytic
continuation of \(E\).  In particular, the unit circle is a natural boundary,
and \(E\) is transcendental over \(\C(z)\).
\end{theorem}

\begin{proof}
Let \(\zeta\) be a dyadic root of unity, say \(\zeta^{2^r}=1\).  For
\(0<\rho<1\), the factors with \(j\ge r\) are real numbers in \([0,1]\), and
\[
 |E(\rho\zeta)|
 \le 2^r(1-\rho^{2^r}).
\]
Hence \(E(\rho\zeta)\to0\) radially as \(\rho\uparrow1\).  Dyadic roots of
unity are dense on the circle.  If \(E\) continued through a nonempty arc,
the continuation would vanish at infinitely many dyadic roots with an
accumulation point in its domain.  The identity theorem would force it to
vanish identically, contradicting \(E(0)=1\).  Thus no arc permits
continuation.  An algebraic function over \(\C(z)\) has only finitely many
branch points and poles on the Riemann sphere, so it cannot have a natural
boundary.
\end{proof}

\begin{remark}
The general rational-or-natural-boundary dichotomy for Mahler functions is due
to Rand\'e and has modern proofs by Bell, Coons, and Rowland
\cite{BellCoonsRowland}.  Here the dense radial-zero argument is specific and
self-contained.
\end{remark}

\section{The ruler sequence, coefficient recurrences, partitions, and determinants}
\label{sec:ruler}

Define the odd ruler values
\begin{equation}
 a_r=2^{\vtwo(r)+1}-1\qquad(r\ge1),
 \label{eq:ruler-a}
\end{equation}
so \(a_1,a_2,\ldots=1,3,1,7,1,3,1,15,\ldots\).  Logarithmic expansion of
\eqref{eq:E-product} gives
\begin{equation}
 \boxed{\log E(z)=-\sum_{r\ge1}\frac{a_r}{r}z^r,}
 \qquad
 -\frac{zE'(z)}{E(z)}=\sum_{r\ge1}a_rz^r.
 \label{eq:log-E-ruler}
\end{equation}
Indeed, the coefficient at \(z^r\) receives \(-2^j/r\) from every
\(2^j\mid r\), whose sum is \(-a_r/r\).

Comparing coefficients in \(-zE'=E\sum a_rz^r\) yields a nonbinary
convolution recurrence:
\begin{equation}
 \boxed{n\eps_n=-\sum_{r=1}^na_r\eps_{n-r}\qquad(n\ge1).}
 \label{eq:ruler-convolution}
\end{equation}
The ruler sequence itself obeys
\begin{equation}
 a_{2n+1}=1,\qquad a_{2n}=2a_n+1.
 \label{eq:ruler-recursion}
\end{equation}
Its Lambert and Dirichlet series are
\begin{align}
 \sum_{r\ge1}a_rz^r
 &=\sum_{j\ge0}\frac{2^jz^{2^j}}{1-z^{2^j}},
 \label{eq:ruler-Lambert}\\
 \sum_{r\ge1}\frac{a_r}{r^s}
 &=\boxed{\frac{\zeta(s)}{1-2^{1-s}}}\qquad(\Re s>1).
 \label{eq:ruler-Dirichlet}
\end{align}
For the second formula, split \(r=2^ku\) with \(u\) odd and sum a geometric
series.

\subsection{Partition and Bell-polynomial formulas}

Exponentiating \eqref{eq:log-E-ruler} gives the finite partition sum
\begin{equation}
 \boxed{\eps_n=
 \sum_{\substack{k_1,\ldots,k_n\ge0\\
                  k_1+2k_2+\cdots+nk_n=n}}
 \prod_{r=1}^n\frac1{k_r!}
 \left(-\frac{a_r}{r}\right)^{k_r}.}
 \label{eq:partition-formula}
\end{equation}
The rational summands cancel to exactly \(\pm1\).  In terms of complete
exponential Bell polynomials,
\begin{equation}
 \boxed{\eps_n=\frac1{n!}\Bell_n
 \bigl(-0!a_1,-1!a_2,\ldots,-(n-1)!a_n\bigr).}
 \label{eq:Bell-formula}
\end{equation}

\subsection{A factorial determinant}

For \(n\ge1\), let
\begin{equation}
 H_n=\begin{pmatrix}
 a_1&1&0&\cdots&0\\
 a_2&a_1&2&\ddots&\vdots\\
 a_3&a_2&a_1&\ddots&0\\
 \vdots&\vdots&\ddots&\ddots&n-1\\
 a_n&a_{n-1}&\cdots&a_2&a_1
 \end{pmatrix}.
 \label{eq:Hessenberg}
\end{equation}
Then
\begin{equation}
 \boxed{\eps_n=\frac{(-1)^n}{n!}\det H_n,\qquad
 |\det H_n|=n!.}
 \label{eq:Hessenberg-formula}
\end{equation}
To verify it, expand a lower-Hessenberg determinant to obtain
\[
 D_n=\sum_{r=1}^n(-1)^{r+1}\frac{(n-1)!}{(n-r)!}a_rD_{n-r},
 \qquad D_0=1,
\]
and set \(c_n=(-1)^nD_n/n!\); then \(c_n\) satisfies
\eqref{eq:ruler-convolution} with \(c_0=1\).

\section{Prouhet cancellation and a master finite-difference identity}
\label{sec:prouhet}

\subsection{Exact vanishing order}

Every factor in \eqref{eq:finite-product} has a simple zero at \(z=1\), so
\begin{equation}
 P_m(z)=(-1)^m2^{\binom m2}(z-1)^m+\bigO((z-1)^{m+1}).
 \label{eq:Pm-local}
\end{equation}
Thus \(P_m\) has a zero of exact order \(m\) at one.

\begin{theorem}[Prouhet cancellation]
For every polynomial \(p\) of degree less than \(m\),
\begin{equation}
 \boxed{\sum_{n=0}^{2^m-1}\eps_np(n)=0.}
 \label{eq:prouhet}
\end{equation}
Equivalently, the evil and odious subsets of \(\{0,\ldots,2^m-1\}\) have
equal sums of like powers through degree \(m-1\).
\end{theorem}

\begin{proof}
The derivative \(P_m^{(r)}(1)\) equals
\(\sum_{n<2^m}\eps_n\fall{n}{r}\) and vanishes for \(r<m\).
Falling factorials form a basis of the polynomial ring.
\end{proof}

The first surviving moments are
\begin{align}
 \boxed{\sum_{n<2^m}\eps_nn^m
   =(-1)^m m!2^{\binom m2},}
 \label{eq:first-power-moment}\\
 \boxed{\sum_{n<2^m}\eps_n\binom nm
   =(-1)^m2^{\binom m2}.}
 \label{eq:first-binomial-moment}
\end{align}

\subsection{The operator formula}

For \(h\in\C\), let \(T_hf(x)=f(x+h)\) and
\(\Delta_h=T_h-I\).  Expanding a product of commuting shifts proves the
master identity
\begin{equation}
 \boxed{\sum_{n=0}^{2^m-1}\eps_nf(x+nh)
 =\left[\prod_{j=0}^{m-1}(I-T_{2^jh})f\right](x)
 =(-1)^m\Delta_h\Delta_{2h}\cdots\Delta_{2^{m-1}h}f(x).}
 \label{eq:finite-difference-master}
\end{equation}
This single equation generates Prouhet cancellation, exact remainders,
reciprocal-power positivity, and sparse Pascal analogues.

If \(h\ge0\) and \(f\in C^m\), repeated use of the fundamental theorem of
calculus gives
\begin{align}
 \sum_{n<2^m}\eps_nf(x+nh)
 &=(-1)^m\int_0^h\int_0^{2h}\cdots\int_0^{2^{m-1}h}
 f^{(m)}(x+t_0+\cdots+t_{m-1})
 \dd t_{m-1}\cdots\dd t_0,
 \label{eq:cubical-integral}\\
 \left|\sum_{n<2^m}\eps_nf(x+nh)\right|
 &\le 2^{\binom m2}h^m
 \sup|f^{(m)}|.
 \label{eq:cubical-bound}
\end{align}

\begin{corollary}[Positive reciprocal-power sums]
For \(a,s>0\),
\begin{equation}
 \boxed{\sum_{n=0}^{2^m-1}\frac{\eps_n}{(a+n)^s}
 =(s)_m\int_0^1\int_0^2\cdots\int_0^{2^{m-1}}
 \frac{\dd t_{m-1}\cdots\dd t_0}
 {(a+t_0+\cdots+t_{m-1})^{s+m}}>0,}
 \label{eq:positive-reciprocal}
\end{equation}
where \((s)_m=s(s+1)\cdots(s+m-1)\).
\end{corollary}

\begin{proof}
Use \eqref{eq:cubical-integral} with \(f(x)=(a+x)^{-s}\).  Its \(m\)-th
derivative is \((-1)^m(s)_m(a+x)^{-s-m}\), so the signs cancel.
\end{proof}

\subsection{All higher moments}

Set
\begin{equation}
 A_m=\frac{2^m-1}{2},\qquad
 c_{m,r}=\frac{B_{2r}}{2r(2r)!}
 \frac{2^{2rm}-1}{2^{2r}-1}\quad(r\ge1),
 \label{eq:moment-constants}
\end{equation}
where \(B_{2r}\) are Bernoulli numbers with \(B_2=1/6\).  Substituting
\(z=e^t\) in the finite product and using the expansion of
\(\log(\sinh x/x)\) gives
\begin{equation}
 \boxed{\sum_{n<2^m}\eps_ne^{nt}
 =(-1)^m2^{\binom m2}t^m
 \exp\!\left(A_mt+\sum_{r\ge1}c_{m,r}t^{2r}\right).}
 \label{eq:exponential-moment-factor}
\end{equation}
Consequently, for every \(\ell\ge0\),
\begin{equation}
 \boxed{\sum_{n<2^m}\eps_nn^{m+\ell}
 =(-1)^m2^{\binom m2}(m+\ell)!
 [t^\ell]\exp\!\left(A_mt+\sum_{r\ge1}c_{m,r}t^{2r}\right).}
 \label{eq:all-power-moments}
\end{equation}
After centering at \(A_m\), only even corrections survive:
\begin{equation}
 \sum_{n<2^m}\eps_ne^{(n-A_m)t}
 =(-1)^m2^{\binom m2}t^m
 \exp\!\left(\sum_{r\ge1}c_{m,r}t^{2r}\right).
 \label{eq:centered-mgf}
\end{equation}
Hence
\begin{equation}
 \sum_{n<2^m}\eps_n(n-A_m)^q=0
 \quad\text{unless }q\ge m\text{ and }q\equiv m\pmod2.
 \label{eq:centered-parity}
\end{equation}
This also follows from the reciprocity \eqref{eq:Pm-reciprocity}.

For binomial moments, the exact coefficient generator is
\begin{equation}
 \boxed{\sum_{r\ge0}\left(\sum_{n<2^m}\eps_n\binom nr\right)u^r
 =\prod_{j=0}^{m-1}\left(1-(1+u)^{2^j}\right).}
 \label{eq:all-binomial-moments}
\end{equation}
It displays the zeros for \(r<m\) and the first value
\eqref{eq:first-binomial-moment} without conversion between power bases.

\section{Sparse Prouhet identities on arbitrary Pascal supports}
\label{sec:sparse}

Let
\begin{equation}
 J(n)=\{j\ge0:b_j(n)=1\},\qquad
 s=|J(n)|=\wt(n),\qquad
 \beta(n)=\sum_{j\in J(n)}j.
 \label{eq:J-beta}
\end{equation}

\begin{theorem}[Pascal-support product]
For every \(n\ge0\),
\begin{equation}
 \boxed{R_n(z):=\sum_{k\submask n}\eps_kz^k
 =\prod_{j\in J(n)}(1-z^{2^j}).}
 \label{eq:sparse-product}
\end{equation}
Equivalently,
\begin{equation}
 R_n(z)=\sum_{k=0}^n\eps_k
 \ind_{\binom nk\text{ odd}}z^k.
 \label{eq:sparse-Pascal}
\end{equation}
\end{theorem}

\begin{proof}
Expand the product.  Each chosen subset of \(J(n)\) gives one submask \(k\)
and coefficient \((-1)^{\wt(k)}\).  Lucas's criterion gives the second form.
\end{proof}

The polynomial \(R_n\) has a zero of order \(s\) at one.  Therefore
\begin{equation}
 \boxed{\sum_{k\submask n}\eps_kp(k)=0\qquad(\deg p<s),}
 \label{eq:sparse-prouhet}
\end{equation}
and the first surviving moment is
\begin{equation}
 \boxed{\sum_{k\submask n}\eps_kk^s
 =(-1)^ss!2^{\beta(n)}.}
 \label{eq:sparse-first-moment}
\end{equation}
In particular, every nonzero Pascal row has equally many evil and odious
indices among the positions of its odd entries.

The operator form is
\begin{equation}
 \boxed{\sum_{k\submask n}\eps_kf(x+hk)
 =\prod_{j\in J(n)}(I-T_{2^jh})f(x).}
 \label{eq:sparse-operator}
\end{equation}
For \(f\in C^s\) and \(h\ge0\), its absolute value is at most
\begin{equation}
 2^{\beta(n)}h^s\sup_{x\le y\le x+nh}|f^{(s)}(y)|.
 \label{eq:sparse-bound}
\end{equation}

\begin{example}
For \(n=13=1101_2\), the odd positions of Pascal row 13 are
\[
 0,1,4,5,8,9,12,13.
\]
Their signed sums against \(1,k,k^2\) vanish, while
\[
 \sum_{k\submask13}\eps_kk^3=-3!\,2^{0+2+3}=-192.
\]
\end{example}

\section{The bit series in characteristic two}
\label{sec:F2}

In \(\F_2[[x]]\), put
\begin{equation}
 \Theta(x)=\sum_{n\ge0}\TM(n)x^n.
 \label{eq:Theta-F2}
\end{equation}
Splitting even and odd coefficients gives
\begin{equation}
 \Theta(x)=(1+x)\Theta(x^2)+\frac{x}{1+x^2}.
 \label{eq:F2-functional}
\end{equation}
Frobenius says \(\Theta(x^2)=\Theta(x)^2\), and
\(1+x^2=(1+x)^2\).  Hence
\begin{equation}
 \boxed{(1+x)^3\Theta(x)^2+(1+x)^2\Theta(x)+x=0.}
 \label{eq:F2-quadratic}
\end{equation}
With \(Y(x)=(1+x)\Theta(x)\), this is the Artin--Schreier equation
\begin{equation}
 \boxed{Y(x)^2+Y(x)=\frac{x}{1+x}.}
 \label{eq:Artin-Schreier}
\end{equation}
Thus the same sequence whose complex generating function is transcendental is
algebraic of degree two over \(\F_2(x)\), an instructive contrast between
characteristic zero and characteristic two.

\section{Iterated prefix sums and terminal zero runs}
\label{sec:iterated-prefix}

Define inclusive discrete antiderivatives by
\begin{equation}
 S_n^{(0)}=\eps_n,
 \qquad
 S_n^{(k+1)}=\sum_{j=0}^nS_j^{(k)}.
 \label{eq:iterated-prefix-def}
\end{equation}
The first one is explicit:
\begin{equation}
 S_{2q}^{(1)}=\eps_q,
 \qquad
 S_{2q+1}^{(1)}=0.
 \label{eq:first-prefix}
\end{equation}
Repeated summation gives
\begin{align}
 \boxed{S_n^{(k)}=
 \sum_{j=0}^n\binom{n-j+k-1}{k-1}\eps_j,}
 \label{eq:prefix-binomial-convolution}\\
 \boxed{\sum_{n\ge0}S_n^{(k)}z^n=\frac{E(z)}{(1-z)^k}.}
 \label{eq:prefix-generating}
\end{align}

\begin{theorem}[Dyadic terminal zero run]
If \(1\le k\le r\), then
\begin{equation}
 S_{2^r-k+i}^{(k)}=0\qquad(0\le i<k),
 \label{eq:terminal-zero-run}
\end{equation}
while
\begin{equation}
 S_{2^r-k-1}^{(k)}=(-1)^{r-k},
 \qquad
 S_{2^r}^{(k)}=-1.
 \label{eq:prefix-boundaries}
\end{equation}
\end{theorem}

\begin{proof}
Below degree \(2^r\), replace \(E(z)\) in
\eqref{eq:prefix-generating} by \(P_r(z)\).  Since
\[
 \frac{P_r(z)}{(1-z)^k}
 =(1-z)^{r-k}\prod_{j=0}^{r-1}(1+z+\cdots+z^{2^j-1})
\]
is a polynomial of degree \(2^r-k-1\), the next \(k\) coefficients vanish.
Its leading coefficient gives the first boundary value; the first omitted
factor \(1-z^{2^r}\) gives the value at \(2^r\).
\end{proof}

\section{Walsh, ordinary Fourier, and cyclotomic formulas}
\label{sec:fourier}

\subsection{The Boolean-cube character}

For \(0\le n<2^m\), let \(\bm b(n)\in\F_2^m\) be its digit vector and let
\(\bm1=(1,\ldots,1)\).  Then
\begin{equation}
 \eps_n=(-1)^{\bm1\cdot\bm b(n)}.
 \label{eq:walsh-character}
\end{equation}
Equivalently, in natural binary order,
\begin{equation}
 (\eps_0,\ldots,\eps_{2^m-1})^{\mathsf T}
 =\binom{1}{-1}^{\!\otimes m}.
 \label{eq:kronecker}
\end{equation}
Thus the finite block is one Walsh--Hadamard character.  Orthogonality gives
\begin{equation}
 \boxed{\sum_{n<2^m}\eps_n(-1)^{\bm a\cdot\bm b(n)}
 =\begin{cases}2^m,&\bm a=\bm1,\\0,&\bm a\ne\bm1.
 \end{cases}}
 \label{eq:walsh-transform}
\end{equation}
The Walsh spectrum is a point mass even though the ordinary Fourier spectrum
is spread over all odd frequencies.

For \(x\in[0,1)\), set
\(r_j(x)=(-1)^{\lfloor2^{j+1}x\rfloor}\).  If
\(n/2^m\le x<(n+1)/2^m\), then
\begin{equation}
 \eps_n=\prod_{j=0}^{m-1}r_j(x).
 \label{eq:rademacher-product}
\end{equation}

\subsection{Trigonometric factorization}

Putting \(z=e^{2\pi\ii x}\) in \eqref{eq:finite-product} gives
\begin{equation}
 \boxed{P_m(e^{2\pi\ii x})
 =(-2\ii)^m e^{\pi\ii(2^m-1)x}
 \prod_{j=0}^{m-1}\sin(\pi2^jx).}
 \label{eq:trig-factorization}
\end{equation}
In particular,
\begin{equation}
 |P_m(e^{2\pi\ii x})|
 =2^m\prod_{j=0}^{m-1}|\sin(\pi2^jx)|.
 \label{eq:trig-modulus}
\end{equation}
This is the finite Fourier counterpart of the sinc products that occur in the
Fourier transform of the Rvachev up-function.

\subsection{Exact discrete Fourier transform}

For integer \(k\), define
\begin{equation}
 \widehat\eps_m(k)=\sum_{n=0}^{2^m-1}\eps_n
 e^{-2\pi\ii kn/2^m}.
 \label{eq:DFT-def}
\end{equation}
Then
\begin{equation}
 \boxed{\widehat\eps_m(k)=
 \prod_{j=0}^{m-1}\left(1-e^{-2\pi\ii k2^j/2^m}\right).}
 \label{eq:DFT-product}
\end{equation}
Every even mode vanishes.  If \(k\) is odd, no factor vanishes and
\begin{equation}
 \widehat\eps_m(k)
 =(2\ii)^m e^{-\pi\ii k(1-2^{-m})}
 \prod_{r=1}^{m}\sin\!\left(\frac{\pi k}{2^r}\right).
 \label{eq:DFT-odd}
\end{equation}
Finite Fourier inversion therefore gives a cyclotomic formula for each term:
\begin{equation}
 \boxed{\eps_n=2^{-m}\sum_{\substack{0\le k<2^m\\k\text{ odd}}}
 \widehat\eps_m(k)e^{2\pi\ii kn/2^m}}
 \qquad(0\le n<2^m).
 \label{eq:DFT-inversion}
\end{equation}
The formula is independent of how large \(m\) is, once \(2^m>n\).

\section{Riesz products, autocorrelation, and spectral type}
\label{sec:spectral}

Define
\begin{equation}
 g_m(x)=2^{-m}|P_m(e^{2\pi\ii x})|^2.
 \label{eq:gm}
\end{equation}
Then
\begin{equation}
 \boxed{g_m(x)=\prod_{j=0}^{m-1}
 \bigl(1-\cos(2\pi2^jx)\bigr).}
 \label{eq:Riesz-density}
\end{equation}
Parseval's identity gives \(\int_0^1g_m(x)\dd x=1\), so
\(g_m(x)\dd x\) is a probability measure.  Its weak limit is the classical
Thue--Morse Riesz-product measure \(\mu_{\rm TM}\).

For dyadic averages, put
\begin{equation}
 \eta_m(r)=2^{-m}\sum_{n=0}^{2^m-1}\eps_n\eps_{n+r}.
 \label{eq:finite-correlation}
\end{equation}
Even--odd splitting gives the exact recurrences
\begin{align}
 \eta_{m+1}(2r)&=\eta_m(r),
 \label{eq:eta-finite-even}\\
 \eta_{m+1}(2r+1)&=-\frac{\eta_m(r)+\eta_m(r+1)}2.
 \label{eq:eta-finite-odd}
\end{align}
For each fixed \(r\), the limit \(\eta(r)=\lim_m\eta_m(r)\) exists and is
characterized by
\begin{equation}
 \boxed{\eta(0)=1,\qquad
 \eta(2r)=\eta(r),\qquad
 \eta(2r+1)=-\frac{\eta(r)+\eta(r+1)}2.}
 \label{eq:eta-limit}
\end{equation}
For example,
\begin{equation}
 \eta(1)=\eta(2)=-\frac13,\qquad
 \eta(3)=\frac13,\qquad \eta(5)=0.
 \label{eq:eta-examples}
\end{equation}
These are the Fourier coefficients of \(\mu_{\rm TM}\).

\begin{literature}[Purely singular continuous spectrum]
The Riesz-product measure \(\mu_{\rm TM}\) has no atoms and is singular with
respect to Lebesgue measure.  Accordingly, the balanced Thue--Morse
substitution has a purely singular continuous spectral component associated
with the weight \((1,-1)\); see Queff\'elec \cite{QueffelecSubstitution} and
Baake--Grimm \cite{BaakeGrimmTM}.  The recursion
\eqref{eq:eta-limit}, the Riesz product \eqref{eq:Riesz-density}, and the
substitution dynamical system are three exact descriptions of the same
measure.
\end{literature}

\begin{literature}[Gelfond's exponential-sum bound]
There is an absolute constant \(C\) such that
\begin{equation}
 \sup_{x\in\R}|P_m(e^{2\pi\ii x})|
 \le C\,3^{m/2}.
 \label{eq:Gelfond-bound}
\end{equation}
Equivalently, for arbitrary \(N\),
\begin{equation}
 \sup_x\left|\sum_{n<N}\eps_ne^{2\pi\ii nx}\right|
 \ll N^{\log3/\log4}.
 \label{eq:Gelfond-exponent}
\end{equation}
The exponent \(\log3/\log4\) is sharp.  This estimate, originating in
Gelfond's work \cite{Gelfond}, is the analytic input behind many distribution
results for digit sums.
\end{literature}

\section{Base-\texorpdfstring{\(b\)}{b} Thue--Morse constants and Mahler transcendence}
\label{sec:constants}

For an integer \(b\ge2\), define
\begin{equation}
 \Theta_b=\sum_{n\ge0}\frac{\TM(n)}{b^{n+1}}.
 \label{eq:Theta-b}
\end{equation}
Substituting \(z=1/b\) in \eqref{eq:T-from-E} gives
\begin{equation}
 \boxed{\Theta_b=
 \frac1{2(b-1)}-\frac1{2b}
 \prod_{j\ge0}(1-b^{-2^j}).}
 \label{eq:Theta-b-product}
\end{equation}
The signed companion is
\begin{equation}
 \sum_{n\ge0}\frac{\eps_n}{b^{n+1}}
 =\frac1b\prod_{j\ge0}(1-b^{-2^j}).
 \label{eq:signed-b-constant}
\end{equation}
For \(b=2\), \(\Theta_2=0.011010011001\ldots{}_2\) is the classical
Thue--Morse--Mahler constant.

Nonperiodicity already proves that \(\Theta_b\) is irrational: a rational
number has an eventually periodic base-\(b\) expansion, and the digits here
are not eventually periodic.  A much deeper statement is known.

\begin{literature}[Mahler]
For every integer \(b\ge2\), the number \(\Theta_b\) is transcendental.
Mahler established this through the functional equation
\(E(z)=(1-z)E(z^2)\); later proofs use automatic sequences, Ridout's theorem,
or modern Mahler-method machinery.  See Mahler \cite{Mahler1929}, Allouche--
Shallit \cite{AlloucheShallitBook}, and Bugeaud--Queff\'elec
\cite{BugeaudQueffelec}.
\end{literature}

\section{Shifted Dirichlet series and Mellin transforms}
\label{sec:dirichlet}

For \(a>0\), initially in \(\Re s>0\), define
\begin{equation}
 D_a(s)=\sum_{n\ge0}\frac{\eps_n}{(n+a)^s}.
 \label{eq:Da-def}
\end{equation}
The convergence follows from the bounded prefix sums in
\cref{sec:balance} and Dirichlet summation.

\begin{theorem}[Entire continuation]
\newmaterial For every \(a>0\), \(D_a(s)\) extends to an entire function of
\(s\).  It satisfies
\begin{equation}
 \boxed{D_a(s)=2^{-s}
 \left(D_{a/2}(s)-D_{(a+1)/2}(s)\right).}
 \label{eq:Da-functional}
\end{equation}
\end{theorem}

\begin{proof}
The functional equation follows by separating even and odd indices.  For the
continuation, group the series into blocks of length \(2^m\):
\begin{equation}
 D_a(s)=2^{-ms}\sum_{q\ge0}\eps_q
 \sum_{r=0}^{2^m-1}\eps_r
 \left(q+\frac{a+r}{2^m}\right)^{-s}.
 \label{eq:Da-blocked}
\end{equation}
The inner sum is an \(m\)-fold mixed difference of \(x^{-s}\); uniformly for
\(s\) in a compact set it is \(\bigO(q^{-\Re s-m})\).  Given any compact
subset of the \(s\)-plane, choose \(m\) large enough to make the outer series
normally convergent there.  The representations agree where the original
series converges, so they patch to an entire function.
\end{proof}

There is also a Mellin representation:
\begin{equation}
 \boxed{D_a(s)=\frac1{\Gamma(s)}\int_0^\infty
 t^{s-1}e^{-at}E(e^{-t})\dd t.}
 \label{eq:Da-Mellin}
\end{equation}
Initially this follows by termwise integration.  Near zero, for every
\(m\ge1\),
\begin{equation}
 0<E(e^{-t})\le 2^{\binom m2}t^m,
 \label{eq:E-small-t-flat}
\end{equation}
using the first \(m\) factors of the product.  This super-polynomial flatness
also explains the entire continuation in \eqref{eq:Da-Mellin}.

At \(s=1\), write
\begin{equation}
 H(a)=D_a(1)=\sum_{n\ge0}\frac{\eps_n}{n+a}.
 \label{eq:H-def}
\end{equation}
Then
\begin{equation}
 \boxed{H(a)=\int_0^1x^{a-1}E(x)\dd x>0,}
 \label{eq:H-integral}
\end{equation}
and
\begin{equation}
 H(a)=\frac12H(a/2)-\frac12H((a+1)/2).
 \label{eq:H-functional}
\end{equation}
The positivity is nonobvious in the alternating series but immediate in the
integral because every factor of \(E(x)\) is positive for \(0<x<1\).

\section{A two-parameter product calculus}
\label{sec:infinite-products}

For \(a,b>0\), define
\begin{equation}
 G(a,b)=\prod_{n\ge0}
 \left(\frac{n+a}{n+b}\right)^{\eps_n}.
 \label{eq:G-def}
\end{equation}
The logarithm converges by summation by parts: the partial sums of \(\eps_n\)
are bounded and the first difference of
\(\log((n+a)/(n+b))\) is summable.  Directly,
\begin{equation}
 G(a,a)=1,\qquad G(a,b)G(b,c)=G(a,c),\qquad
 G(b,a)=G(a,b)^{-1}.
 \label{eq:G-cocycle}
\end{equation}
Pairing even and odd indices yields
\begin{equation}
 \boxed{G(a,b)=
 \frac{G(a/2,b/2)}{G((a+1)/2,(b+1)/2)}.}
 \label{eq:G-functional}
\end{equation}

\subsection{The Woods--Robbins product}

The most famous evaluation is
\begin{equation}
 \boxed{G\left(\frac12,1\right)=
 \prod_{n\ge0}\left(\frac{2n+1}{2n+2}\right)^{\eps_n}
 =\frac1{\sqrt2}.}
 \label{eq:Woods-Robbins}
\end{equation}
For a short proof, put
\[
 A=\prod_{n\ge0}\left(\frac{2n+1}{2n+2}\right)^{\eps_n},
 \qquad
 B=\prod_{n\ge1}\left(\frac{2n}{2n+1}\right)^{\eps_n}.
\]
Splitting the product \(AB\) into even and odd indices gives
\(AB=\tfrac12A^{-1}B\), hence \(A^2=1/2\), and positivity selects the
positive root.

A second closed value, proved in the modern product calculus, is
\begin{equation}
 \boxed{G\left(\frac14,\frac34\right)=\frac12.}
 \label{eq:quarter-product}
\end{equation}
See Riasat \cite{RiasatProducts} and Allouche--Riasat--Shallit
\cite{AlloucheRiasatShallit} for systematic extensions.

\subsection{Three dyadic block-product limits}

The product evaluations imply
\begin{align}
 \lim_{m\to\infty}\prod_{k=0}^{2^m-1}
 \left(k+\frac12\right)^{\eps_k}&=\frac12,
 \label{eq:block-product-half}\\
 \lim_{m\to\infty}\prod_{k=0}^{2^m-1}(k+1)^{\eps_k}
 &=\frac1{\sqrt2},
 \label{eq:block-product-one}\\
 \lim_{m\to\infty}\prod_{k=0}^{2^m-1}
 (k+1)^{(-1)^k\eps_k}&=\frac1{2\sqrt2}.
 \label{eq:block-product-alternating}
\end{align}
For example, pairing adjacent factors in the first two finite products gives
\begin{align*}
 \prod_{k<2^m}\left(k+\frac12\right)^{\eps_k}
 &=\prod_{j<2^{m-1}}
 \left(\frac{j+1/4}{j+3/4}\right)^{\eps_j},\\
 \prod_{k<2^m}(k+1)^{\eps_k}
 &=\prod_{j<2^{m-1}}
 \left(\frac{j+1/2}{j+1}\right)^{\eps_j}.
\end{align*}
The third limit is the product of the first two.  Dyadic truncation is natural
because it preserves the exact self-similar blocks.

\section{Combinatorics on the infinite word}
\label{sec:words}

\subsection{Uniform recurrence and reversal symmetry}

The substitution matrix of \(\mu(0)=01,\mu(1)=10\) is strictly positive:
each image contains both letters.  More concretely, every factor occurring in
the fixed point occurs in some finite word \(\mu^k(0)\); every sufficiently
large supertile contains copies of both \(\mu^k(0)\) and \(\mu^k(1)\).
Therefore every factor recurs with bounded gaps.

\begin{proposition}[Uniform recurrence]
The Thue--Morse word is uniformly recurrent: for every finite factor \(w\),
there exists \(L(w)\) such that every factor of length \(L(w)\) contains an
occurrence of \(w\).
\end{proposition}

The reciprocity \eqref{eq:Pm-reciprocity} gives a precise reversal law for the
finite blocks.  If \(0\le n<2^m\), then
\begin{equation}
 \eps_{2^m-1-n}=(-1)^m\eps_n.
 \label{eq:block-reversal}
\end{equation}
Thus a level-\(m\) signed block is a palindrome when \(m\) is even and the
negative of its reversal when \(m\) is odd.  In zero--one notation, odd levels
are complement-palindromes.

\subsection{An elementary linear complexity bound}

Let \(p(\ell)\) denote the number of distinct factors of length \(\ell\).
Choose \(k\) minimal with \(2^k\ge\ell\).  Every length-\(\ell\) factor lies
inside two consecutive level-\(k\) supertiles
\(\mu^k(a)\mu^k(b)\), where \(a,b\in\{0,1\}\).  There are at most four
choices for \(ab\) and at most \(2^k\) starting offsets.  Hence
\begin{equation}
 \boxed{p(\ell)\le4\cdot2^k<8\ell\qquad(\ell\ge2).}
 \label{eq:complexity-bound}
\end{equation}
Sharper exact piecewise-linear formulas are known, but this elementary bound
already has two consequences: the topological entropy is zero, and the word
is not normal.  Indeed, a normal binary word contains all \(2^\ell\) words of
length \(\ell\), incompatible with linear complexity for large \(\ell\).

\begin{literature}[Thue's overlap theorem]
The infinite Thue--Morse word is overlap-free: it contains no factor of the
form \(axaxa\), where \(a\) is a letter and \(x\) a possibly empty word.  In
particular it contains no cube \(www\).  This was one of Thue's foundational
results in combinatorics on words \cite{Thue1906}; modern treatments appear in
Allouche--Shallit \cite{AlloucheShallitBook}.  The statement is about factors
of the original word and should not be confused with the subsequence
\(\TM(n^3)\) discussed later.
\end{literature}

\section{Distribution in arithmetic and polynomial subsequences}
\label{sec:distribution}

The deterministic word has low factor complexity, yet arithmetically sampled
subsequences can display strong pseudorandomness.  The following theorems are
not reproved here; they locate the exact identities of the atlas inside modern
analytic number theory.

\begin{literature}[Newman's phenomenon]
Let
\begin{equation}
 f_{3,0}(N)=\sum_{0\le n<N}\eps_{3n}.
 \label{eq:Newman-sum}
\end{equation}
Then \(f_{3,0}(N)>0\) for every \(N\ge1\): among the first \(N\) multiples of
three, evil numbers always outnumber odious numbers.  Newman proved the
quantitative bounds
\begin{equation}
 \frac1{20}N^{\log3/\log4}
 \le f_{3,0}(N)\le5N^{\log3/\log4},
 \label{eq:Newman-bounds}
\end{equation}
and Coquet later described the exact fractal fluctuation
\cite{Newman,Coquet,ShallitRarefied}.
\end{literature}

\begin{literature}[Prime indices]
Mauduit and Rivat proved that binary digit-sum parity is equidistributed along
the primes.  In signed form,
\begin{equation}
 \sum_{p\le x}\eps_p=\smallo(\pi(x)),
 \label{eq:primes-balance}
\end{equation}
so asymptotically half the primes are evil and half odious
\cite{MauduitRivatPrimes}.
\end{literature}

\begin{literature}[Squares are normal]
Drmota, Mauduit, and Rivat proved that the subsequence
\((\TM(n^2))_{n\ge0}\) is normal: every binary word of length \(L\) occurs
with limiting frequency \(2^{-L}\) \cite{DrmotaMauduitRivatSquares}.  This is
far stronger than the one-bit balance
\(\sum_{n<N}\eps_{n^2}=\smallo(N)\).
\end{literature}

\begin{literature}[Cubes are simply normal]
Spiegelhofer proved
\begin{equation}
 \#\{n<N:\TM(n^3)=0\}\sim\frac N2,
 \qquad\text{equivalently}\qquad
 \sum_{n<N}\eps_{n^3}=\smallo(N).
 \label{eq:cubes-balance}
\end{equation}
As of the cited result, full block normality along cubes was not asserted;
the theorem establishes one-symbol balance and resolves the cubic parity case
of a longstanding Gelfond problem \cite{SpiegelhoferCubes}.
\end{literature}

\begin{literature}[Level of distribution]
Spiegelhofer proved a Bombieri--Vinogradov-type theorem with level of
distribution one (up to the endpoint) for the Thue--Morse sequence.  Among its
applications, \((\TM(\lfloor n^c\rfloor))\) is simply normal for
\(1<c<2\) \cite{SpiegelhoferLevel}.  This illustrates how the finite Fourier
product and carry control feed into arithmetic-progressions estimates.
\end{literature}

\section{Base-\texorpdfstring{\(B\)}{B} root-of-unity generalization}
\label{sec:baseB}

Let \(B\ge2\), let \(s_B(n)\) be the base-\(B\) digit sum, and let
\(\zeta\ne1\) be a \(B\)-th root of unity.  Define
\begin{equation}
 u_n=\zeta^{s_B(n)}.
 \label{eq:baseB-sequence}
\end{equation}
Then
\begin{equation}
 u_{Bn+d}=\zeta^d u_n\qquad(0\le d<B).
 \label{eq:baseB-recursion}
\end{equation}
The finite generating polynomial is
\begin{equation}
 \boxed{P_{m,B,\zeta}(z)=\sum_{n<B^m}u_nz^n
 =\prod_{j=0}^{m-1}\left(\sum_{d=0}^{B-1}
 \zeta^dz^{dB^j}\right).}
 \label{eq:baseB-product}
\end{equation}
Because \(\zeta^B=1\), each factor can also be written as
\begin{equation}
 \sum_{d=0}^{B-1}\zeta^dz^{dB^j}
 =\frac{1-z^{B^{j+1}}}{1-\zeta z^{B^j}}.
 \label{eq:baseB-rational-factor}
\end{equation}

At \(z=1\), every factor has a simple zero and
\begin{equation}
 \left.\frac{\dd}{\dd z}\sum_{d=0}^{B-1}
 \zeta^dz^{dB^j}\right|_{z=1}
 =-\frac{B^{j+1}}{1-\zeta}.
 \label{eq:baseB-factor-derivative}
\end{equation}
Therefore \(P_{m,B,\zeta}\) has a zero of exact order \(m\), and the
base-\(B\) Prouhet theorem follows.

\begin{theorem}[Root-of-unity Prouhet hierarchy]
For every polynomial \(p\) with \(\deg p<m\),
\begin{equation}
 \boxed{\sum_{n=0}^{B^m-1}\zeta^{s_B(n)}p(n)=0.}
 \label{eq:baseB-Prouhet}
\end{equation}
The first surviving moment is
\begin{equation}
 \boxed{\sum_{n=0}^{B^m-1}\zeta^{s_B(n)}n^m
 =(-1)^m m!\frac{B^{m(m+1)/2}}{(1-\zeta)^m}.}
 \label{eq:baseB-first-moment}
\end{equation}
\end{theorem}

\begin{proof}
Differentiate \eqref{eq:baseB-product} at one as in the binary proof.  The
leading coefficient is the product of
\eqref{eq:baseB-factor-derivative}.  Falling factorials again form a basis.
\end{proof}

Fourier inversion over the cyclic group of \(B\)-th roots shows that the
classes
\begin{equation}
 \{0\le n<B^m:s_B(n)\equiv r\pmod B\},\qquad 0\le r<B,
 \label{eq:baseB-classes}
\end{equation}
have equal sums of like powers through degree \(m-1\).  The binary
Thue--Morse Prouhet partition is exactly \(B=2,\zeta=-1\).

For a prime \(p\), Lucas's theorem similarly gives a Pascal-row extension.  If
\(n=\sum n_jp^j\), then
\begin{equation}
 \#\left\{0\le k\le n:\binom nk\not\equiv0\pmod p\right\}
 =\prod_j(n_j+1).
 \label{eq:Lucas-nonzero-p}
\end{equation}
Thus the binary count \(2^{\wt(n)}\) is the \(p=2\) specialization of a
full digitwise product.

\section{The discrete bridge to the Fabius and up-functions}
\label{sec:fabius}

The appearance of Thue--Morse in the Fabius-function development is governed
by one algebraic pipeline:
\begin{equation}
 \boxed{\text{binary character}
 \ \Longrightarrow\ P_r(z)=\prod_{j<r}(1-z^{2^j})
 \ \Longrightarrow\ \frac{P_r(z)}{(1-z)^k}
 \ \Longrightarrow\ \text{iterated-prefix coefficients}.}
 \label{eq:fabius-pipeline}
\end{equation}
The last arrow is exactly \eqref{eq:prefix-generating}.  The factor
\((1-z)^r\mid P_r(z)\) is the algebraic source of the terminal zero runs and
of the high-order flatness that survives after scaling.

On the Fourier side, a normalized factor
\[
 \frac{1-e^{-\ii h\xi}}{\ii h\xi}
 =e^{-\ii h\xi/2}\sinc(h\xi/2)
\]
turns a finite difference into a sinc factor.  Iterating over dyadic steps
produces finite products converging to
\begin{equation}
 \Phi(\xi)=\prod_{j\ge0}\sinc\!\left(\frac{\pi\xi}{2^j}\right),
 \label{eq:up-Fourier-product}
\end{equation}
up to the conventional indexing and Fourier normalization.  This is the
Fourier transform of a compactly supported infinite convolution of uniform
measures, namely the probability law underlying the Rvachev up-function.
The sign identity \eqref{eq:sinc-sign} then explains, in a particularly
literal way, why Thue--Morse signs are encoded in the same sinc product.

After normalization by \(2^{\binom{k}{2}}\), coefficient grids built from
\(S_n^{(k)}\), together with suitable polygonal or spline interpolation,
converge to the Fabius/up profile in the companion development.  The exact
normalization is already visible in the volume factor in
\eqref{eq:cubical-integral} and in the first surviving moment
\eqref{eq:first-power-moment}.  The Prouhet zero supplies flatness, while the
Fourier product supplies compact support and smoothness.  These are not two
separate coincidences; they are the coefficient-space and frequency-space
faces of the same product \(P_r\).

\begin{remark}
The present atlas records the exact algebraic bridge rather than repeating the
quantitative convergence analysis developed elsewhere in the Fabius-function
project.  The relevant repository modules named in the source drafts are:
\begin{center}
\texttt{ThueMorsePrefix.lean} \qquad and \qquad
\texttt{ThueMorseGenerating.lean}.
\end{center}
\end{remark}

\section{Representation comparison and a Lean roadmap}
\label{sec:formalization}

\begin{longtable}{@{}p{0.23\textwidth}p{0.31\textwidth}p{0.37\textwidth}@{}}
\toprule
Representation & Primitive data & Especially effective for \\
\midrule
\endhead
Binary recursion & parity of appended digit & evaluation, induction, blocks \\
Two-kernel automaton & two states and digit matrices & automaticity, extraction \\
XOR character & Boolean group \(\F_2^m\) & Walsh analysis, Hamming spheres \\
Valuation & \(\vtwo(n!)\), Kummer carries & binomial/Catalan arithmetic \\
Pascal support & Lucas submask criterion & log-free and sparse identities \\
Finite Euler product & independent digit choices & Prouhet, DFT, reciprocity \\
Ruler logarithm & divisors \(2^j\mid n\) & convolution, Bell, determinant \\
Finite difference & commuting translations & remainders, bounds, positivity \\
Riesz product & squared Fourier modulus & correlation and spectral measure \\
Mahler equation & scale map \(z\mapsto z^2\) & constants, transcendence, boundary \\
Prefix quotient & multiplication by \((1-z)^{-k}\) & discrete integration, Fabius \\
\bottomrule
\end{longtable}

A dependency-respecting Lean formalization can proceed in the following order.
\begin{enumerate}
 \item Prove the sign normalization and the two recurrences
       \eqref{eq:eps-recursion}.
 \item Establish block multiplicativity, the exact two-kernel, XOR
       multiplicativity, and prefix discrepancy.
 \item Import or prove Legendre, Lucas, and Kummer specializations; derive the
       central-binomial and Catalan formulas.
 \item Define \(P_m\), prove the finite product, reciprocity, and the order of
       vanishing at one.
 \item Formalize the shift-operator identity, Prouhet cancellation, the first
       surviving moment, and the sparse Pascal-support version.
 \item Develop the logarithmic ruler recurrence; Bell polynomials and the
       Hessenberg determinant can be later corollaries.
 \item Add finite DFT and Walsh transforms over \(\mathbb Z/2^m\mathbb Z\) and
       \(\F_2^m\).
 \item Connect repeated prefix sums to \(E(z)/(1-z)^k\) and prove terminal
       zero runs before tackling analytic convergence to Fabius.
\end{enumerate}
Likely theorem names should distinguish the bit \(\TM\), the sign \(\eps\),
and the weight \(\wt\); many informal sign errors arise from silently moving
between these normalizations.

\appendix

\section{The first 64 terms}
\label{app:first64}

The following table displays \(\TM(n)\).  Replacing \(0\) by \(+1\) and
\(1\) by \(-1\) gives \(\eps_n\).

\begin{center}
\small
\begin{longtable}{@{}r|*{8}{c}@{}}
\toprule
\(n\) & \multicolumn{8}{c}{successive values}\\
\midrule
\endhead
0--7   &0&1&1&0&1&0&0&1\\
8--15  &1&0&0&1&0&1&1&0\\
16--23 &1&0&0&1&0&1&1&0\\
24--31 &0&1&1&0&1&0&0&1\\
32--39 &1&0&0&1&0&1&1&0\\
40--47 &0&1&1&0&1&0&0&1\\
48--55 &0&1&1&0&1&0&0&1\\
56--63 &1&0&0&1&0&1&1&0\\
\bottomrule
\end{longtable}
\end{center}

Each row of length eight is a copy or complement of the first row according
to the sign of its block index, as predicted by
\eqref{eq:block-multiplicativity}.

\section{A compact identity index}
\label{app:index}

\begin{align*}
 \TM(n)&\longleftrightarrow \eps_n=1-2\TM(n),\\
 \eps_n&\longleftrightarrow(-1)^{\wt(n)},\\
 \wt(n)&\longleftrightarrow n-\vtwo(n!),\\
 \wt(n)&\longleftrightarrow\vtwo\binom{2n}{n},\\
 b_j(n)&\longleftrightarrow\binom{n}{2^j}\bmod2,\\
 k\submask n&\longleftrightarrow\binom nk\text{ odd},\\
 \eps_n&\longleftrightarrow[z^n]\prod_{j\ge0}(1-z^{2^j}),\\
 \sum_{n<2^m}\eps_nf(n)
 &\longleftrightarrow(-1)^m\Delta_1\Delta_2\cdots
 \Delta_{2^{m-1}}f(0),\\
 \eta(r)&\longleftrightarrow\widehat\mu_{\rm TM}(r),\\
 S_n^{(k)}&\longleftrightarrow[z^n]\frac{E(z)}{(1-z)^k}.
\end{align*}

\section{Executable checks}
\label{app:code}

The following short Python program checks a representative cross-section of
the identities.  It deliberately uses integer arithmetic except in the
Fourier test.

\begin{lstlisting}[language=Python]
from math import comb, factorial
import cmath


def wt(n: int) -> int:
    return n.bit_count()


def eps(n: int) -> int:
    return -1 if wt(n) & 1 else 1


def v2(n: int) -> int:
    if n <= 0:
        raise ValueError("v2 expects a positive integer")
    return (n & -n).bit_length() - 1


def catalan(n: int) -> int:
    return comb(2*n, n) // (n + 1)

for n in range(1, 500):
    assert v2(comb(2*n, n)) == wt(n)
    assert v2(catalan(n)) == wt(n + 1) - 1

for m in range(1, 9):
    # Prouhet moments.
    for r in range(m):
        assert sum(eps(n) * n**r for n in range(2**m)) == 0
    assert sum(eps(n) * n**m for n in range(2**m)) == \
           (-1)**m * factorial(m) * 2**(m*(m-1)//2)

    # DFT: every even frequency vanishes.
    for k in range(0, 2**m, 2):
        s = sum(eps(n) * cmath.exp(-2j*cmath.pi*k*n/2**m)
                for n in range(2**m))
        assert abs(s) < 1e-8

for n in range(1, 200):
    odd_positions = [k for k in range(n + 1) if comb(n, k) & 1]
    assert len(odd_positions) == 2**wt(n)
    assert sum(eps(k) for k in odd_positions) == 0

print("all checks passed")
\end{lstlisting}

\section{Bibliographic notes}

The Thue--Morse sequence has been rediscovered in number theory,
combinatorics on words, dynamics, harmonic analysis, automata, and fair
allocation.  The bibliography below is selective: it emphasizes sources that
support the formula families merged here and the extra material added in this
version.  The four repository drafts themselves contain further specialized
references.

\begin{thebibliography}{99}

\bibitem{AlloucheShallitUbiquitous}
J.-P. Allouche and J. Shallit,
\emph{The ubiquitous Prouhet--Thue--Morse sequence},
in Sequences and Their Applications, Springer, 1999, pp. 1--16.

\bibitem{AlloucheShallitBook}
J.-P. Allouche and J. Shallit,
\emph{Automatic Sequences: Theory, Applications, Generalizations},
Cambridge University Press, 2003.

\bibitem{AlloucheRiasatShallit}
J.-P. Allouche, S. Riasat, and J. Shallit,
More infinite products involving the Thue--Morse sequence,
\emph{J. Integer Sequences} 22 (2019), Article 19.2.3;
see also arXiv:1709.03398.

\bibitem{BaakeGrimmTM}
M. Baake and U. Grimm,
The singular continuous diffraction measure of the Thue--Morse chain,
\emph{J. Phys. A} 41 (2008), 422001.

\bibitem{BellCoonsRowland}
J. P. Bell, M. Coons, and E. Rowland,
The rational--transcendental dichotomy of Mahler functions,
\emph{J. Integer Sequences} 16 (2013), Article 13.2.10.

\bibitem{BugeaudQueffelec}
Y. Bugeaud and M. Queff\'elec,
On the rational approximation to the binary Thue--Morse--Mahler number,
\emph{J. Integer Sequences} 16 (2013), Article 13.2.1.

\bibitem{Coquet}
J. Coquet,
A summation formula related to the binary digits,
\emph{Invent. Math.} 73 (1983), 107--115.

\bibitem{DrmotaMauduitRivatSquares}
M. Drmota, C. Mauduit, and J. Rivat,
The Thue--Morse sequence along the squares is normal,
preprint dated 2013; published version in the analytic-number-theory
literature.

\bibitem{Gelfond}
A. O. Gelfond,
Sur les nombres qui ont des propri\'et\'es additives et multiplicatives
donn\'ees,
\emph{Acta Arith.} 13 (1967/68), 259--265.

\bibitem{Kummer}
E. E. Kummer,
\"Uber die Erg\"anzungss\"atze zu den allgemeinen Reciprocit\"atsgesetzen,
\emph{J. Reine Angew. Math.} 44 (1852), 93--146.

\bibitem{Lucas}
\'E. Lucas,
Th\'eorie des fonctions num\'eriques simplement p\'eriodiques,
\emph{Amer. J. Math.} 1 (1878), 184--240 and 289--321.

\bibitem{Mahler1929}
K. Mahler,
Arithmetische Eigenschaften der L\"osungen einer Klasse von
Funktionalgleichungen,
\emph{Math. Ann.} 101 (1929), 342--366; corrigendum 103 (1930), 532.

\bibitem{MauduitRivatPrimes}
C. Mauduit and J. Rivat,
Sur un probl\`eme de Gelfond: la somme des chiffres des nombres premiers,
\emph{Ann. of Math.} 171 (2010), 1591--1646.

\bibitem{Morse1921}
M. Morse,
Recurrent geodesics on a surface of negative curvature,
\emph{Trans. Amer. Math. Soc.} 22 (1921), 84--100.

\bibitem{Newman}
D. J. Newman,
On the number of binary digits in a multiple of three,
\emph{Proc. Amer. Math. Soc.} 21 (1969), 719--721.

\bibitem{Prouhet1851}
E. Prouhet,
M\'emoire sur quelques relations entre les puissances des nombres,
\emph{C. R. Acad. Sci. Paris} 33 (1851), 225.

\bibitem{QueffelecSubstitution}
M. Queff\'elec,
\emph{Substitution Dynamical Systems -- Spectral Analysis},
2nd ed., Lecture Notes in Mathematics 1294, Springer, 2010.

\bibitem{RiasatProducts}
S. Riasat,
Infinite products involving binary digit sums,
arXiv:1709.04104, version of 2018.

\bibitem{Robbins}
D. Robbins,
Solution to Problem E2692,
\emph{Amer. Math. Monthly} 86 (1979), 394--395.

\bibitem{ShallitRarefied}
J. Shallit,
Rarefied Thue--Morse sums via automata theory and logic,
arXiv:2302.09436, 2023.

\bibitem{SpiegelhoferCubes}
L. Spiegelhofer,
Thue--Morse along the sequence of cubes,
arXiv:2308.09498, 2023.

\bibitem{SpiegelhoferLevel}
L. Spiegelhofer,
The level of distribution of the Thue--Morse sequence,
\emph{Compos. Math.} 156 (2020), 2560--2587.

\bibitem{Thue1906}
A. Thue,
\"Uber unendliche Zeichenreihen,
\emph{Norske Vid. Selsk. Skr. I. Mat.-Nat. Kl.} 7 (1906), 1--22.

\bibitem{Woods}
D. R. Woods,
Problem E2692,
\emph{Amer. Math. Monthly} 85 (1978), 48.

\bibitem{ReshetnikovRepo}
V. Reshetnikov,
\repo: Analysis/FabiusFunction,
\url{https://github.com/VladimirReshetnikov/ProveIt/tree/main/Analysis/FabiusFunction},
accessed 26 August 2026.

\end{thebibliography}

\end{document}
